\documentclass[tikz,border=4pt]{standalone}
\usepackage{ctex}
\usetikzlibrary{positioning}
\begin{document}
\begin{tikzpicture}[
    box/.style={rectangle, draw, rounded corners=4pt, minimum width=7cm, minimum height=0.6cm, font=\small, align=left, fill=blue!5},
    title/.style={font=\bfseries\large},
]
\node[title] (t) at (0,0) {不定积分基本公式表};
\node[box, below=0.3cm of t] (r1) {$\int x^n\,dx = \frac{x^{n+1}}{n+1}+C$（$n\neq -1$）};
\node[box, below=0.1cm of r1] (r2) {$\int \frac{1}{x}\,dx = \ln|x|+C$};
\node[box, below=0.1cm of r2] (r3) {$\int e^x\,dx = e^x+C$};
\node[box, below=0.1cm of r3] (r4) {$\int a^x\,dx = \frac{a^x}{\ln a}+C$};
\node[box, below=0.1cm of r4] (r5) {$\int \sin x\,dx = -\cos x+C$};
\node[box, below=0.1cm of r5] (r6) {$\int \cos x\,dx = \sin x+C$};
\node[box, below=0.1cm of r6] (r7) {$\int \sec^2 x\,dx = \tan x+C$};
\node[box, below=0.1cm of r7] (r8) {$\int \frac{1}{\sqrt{1-x^2}}\,dx = \arcsin x+C$};
\node[box, below=0.1cm of r8] (r9) {$\int \frac{1}{1+x^2}\,dx = \arctan x+C$};
\node[font=\footnotesize, below=0.3cm of r9] {**$\pm C$ 不可漏**——不定积分是函数族};
\end{tikzpicture}
\end{document}
